\documentclass[11pt,a4paper]{article}

% Packages
\usepackage[utf8]{inputenc}
\usepackage[T1]{fontenc}
\usepackage{lmodern}
\usepackage[margin=2.5cm]{geometry}
\usepackage{hyperref}
\usepackage{graphicx}
\usepackage{booktabs}
\usepackage{longtable}
\usepackage{listings}
\usepackage{xcolor}
\usepackage{float}
\usepackage{caption}
\usepackage{enumitem}
\usepackage{fancyvrb}
\usepackage{pmboxdraw}

% Hyperref setup
\hypersetup{
    colorlinks=true,
    linkcolor=blue,
    urlcolor=blue,
    citecolor=blue,
}

% Code listings style
\lstset{
    basicstyle=\ttfamily\small,
    breaklines=true,
    frame=single,
    backgroundcolor=\color{gray!10},
    showstringspaces=false,
    keywordstyle=\color{blue},
    commentstyle=\color{green!50!black},
    stringstyle=\color{red!60!black},
    language=Python,
    columns=flexible,
    keepspaces=true,
}

% Bash listings style
\lstdefinestyle{bash}{
    language=bash,
    basicstyle=\ttfamily\small,
    breaklines=true,
    frame=single,
    backgroundcolor=\color{gray!10},
    showstringspaces=false,
    columns=flexible,
    keepspaces=true,
}

% Ini listings style
\lstdefinestyle{ini}{
    basicstyle=\ttfamily\small,
    breaklines=true,
    frame=single,
    backgroundcolor=\color{gray!10},
    showstringspaces=false,
    columns=flexible,
    keepspaces=true,
}

% Dockerfile listings style
\lstdefinestyle{docker}{
    basicstyle=\ttfamily\small,
    breaklines=true,
    frame=single,
    backgroundcolor=\color{gray!10},
    showstringspaces=false,
    columns=flexible,
    keepspaces=true,
}

% Title info
\title{\textbf{CISC 886 -- Cloud Computing}\\\Large Project Deliverable Report\\\normalsize School of Computing, Queen's University, Kingston, Canada}
\author{Student: [YOUR\_NAME] \quad Student ID: [YOUR\_STUDENT\_ID] \\ Queen's NetID: [YOUR\_NETID] (e.g., 25tfgf)}
\date{[COMPLETION\_DATE]}

\begin{document}

\maketitle

\section{System Architecture (2 marks)}

\subsection{Architecture Diagram}

\begin{verbatim}
+------------------------------------------------------------------------------+
|                              AWS Cloud - us-east-1                           |
|                                                                              |
|  +========================================================================+  |
|  |                         25tfgf-vpc (10.0.0.0/16)                        |  |
|  |                                                                        |  |
|  |  +----------------------------------------------------------------+   |  |
|  |  |                   Public Subnets (2 AZs)                        |   |  |
|  |  |                                                                |   |  |
|  |  |  +-------------------+        +-------------------+            |   |  |
|  |  |  |  Public Subnet 1   |        |  Public Subnet 2   |            |   |  |
|  |  |  |  10.0.1.0/24       |        |  10.0.2.0/24       |            |   |  |
|  |  |  |  (us-east-1a)      |        |  (us-east-1b)      |            |   |  |
|  |  |  |                   |        |                   |            |   |  |
|  |  |  |  +-------------+  |        |                   |            |   |  |
|  |  |  |  | EMR Master  |  |        |                   |            |   |  |
|  |  |  |  | m5.xlarge   |  |        |                   |            |   |  |
|  |  |  |  +-------------+  |        |                   |            |   |  |
|  |  |  +-------------------+        +-------------------+            |   |  |
|  |  |                                                             |   |  |
|  |  +----------------------------------------------------------------+   |  |
|  |                                                                      |   |
|  +========================================================================+  |
|                                                                              |
+------------------------------------------------------------------------------+

                        DATA FLOW ARCHITECTURE

+-----------+     +---------------+     +-------------+     +------------+
| HuggingFace|     |      S3       |     |     EMR      |     |     S3      |
|  Dataset   |     |(25tfgf-ai-    |     |   (Spark)    |     |(processed) |
| (Download) |--->>|   medical)    |-->  |  m5.xlarge   | --> |   Parquet   |
+-----------+     |   (raw/)      |     |  1M+2C nodes |     |            |
                  +---------------+     +-------------+     +------+-----+
                                                                |
                                                                v
                                                       +--------------+
                                                       |  Local Work- |
                                                       |  station     |
                                                       | (Unsloth     |
                                                       | QLoRA 50k)   |
                                                       +------+-------+
                                                                |
                                                                v
                                                       +--------------+
                                                       |  S3 Bucket   |
                                                       |(upload model)|
                                                       +------+-------+
                                                                |
                                                                v
+============================================================================+
||                           25tfgf-vpc                                 ||
||  +---------------------------------------------------------------+       ||
||  |  Security Groups:                                             |       ||
||  |  - 25tfgf-ec2-sg: SSH(22), Ollama(11434), OpenWebUI(8080)     |       ||
||  |  - 25tfgf-emr-master-sg: SSH(22), JupyterHub(9443), Internal ||       ||
||  |  - 25tfgf-emr-slave-sg: SSH(22), Internal                    ||       ||
||  +---------------------------------------------------------------+       ||
||                                                                       ||
||                         +------------------+                         ||
||                         |  EC2 m5.xlarge   |                         ||
||                         |  Ubuntu 22.04    |                         ||
||                         |                  |                         ||
||                         |  +-----------+   |                         ||
||                         |  |  Ollama   |   |                         ||
||                         |  |  :11434   |   |                         ||
||                         |  +-----------+   |                         ||
||                         |  +-----------+   |                         ||
||                         |  | OpenWebUI |   |                         ||
||                         |  |   :8080   |   |                         ||
||                         |  +-----------+   |                         ||
||                         +---------+---------+                         ||
+============================================================================+
                                        |
                                        | Port 8080 (OpenWebUI)
                                        v
                               +------------------+
                               |   User Browser   |
                               |  (Chat Interface) |
                               +------------------+

LEGEND:
  ---> : Data flow direction
  1M+2C : 1 Master + 2 Core nodes
  QLoRA : Quantized Low-Rank Adaptation
  AZ    : Availability Zone
\end{verbatim}

\subsection{Data Flow Description}

\begin{quote}
Data flows from the HuggingFace Medical Dataset (\texttt{ruslanmv/ai-medical-dataset}) downloaded locally and uploaded to S3 as raw data in the \texttt{25tfgf-ai-medical} bucket. The EMR Spark cluster (1 master + 2 core nodes, m5.xlarge) processes this data through a PySpark pipeline that applies length-based quality filtering (context: 200--4096 chars, question: $\geq$20 chars), creates a medical prompt template, and splits into train/val/test sets (80/10/10) stored as Parquet in S3. The processed data is downloaded to a local workstation (RTX 5000 Ada, 16 GB VRAM) where QLoRA fine-tuning is performed on \texttt{unsloth/Llama-3.2-1B-Instruct} using Unsloth. The fine-tuned LoRA adapter is saved and deployed to an EC2 m5.xlarge instance in the 25tfgf-vpc, where Ollama serves the model at port 11434 and OpenWebUI provides a browser-based chat interface at port 8080. The auto-start systemd services ensure both Ollama and OpenWebUI restart automatically after any server reboot.
\end{quote}

\section{VPC \& Networking (4 marks)}

\subsection{VPC Configuration}

\begin{table}[H]
\centering
\begin{tabular}{@{}ll@{}}
\toprule
\textbf{Component} & \textbf{Value} \\
\midrule
VPC Name & 25tfgf-vpc \\
CIDR Block & 10.0.0.0/16 \\
Region & us-east-1 \\
Number of AZs & 2 (us-east-1a, us-east-1b) \\
\bottomrule
\end{tabular}
\end{table}

\textbf{Justification for CIDR block:} The /16 provides 65,536 IP addresses which is sufficient for this project while leaving room for expansion. Using a private CIDR block (10.x.x.x) ensures the VPC is not directly exposed to the internet.

\subsection{Subnet Design}

\begin{table}[H]
\centering
\begin{tabular}{@{}llll@{}}
\toprule
\textbf{Subnet} & \textbf{CIDR} & \textbf{AZ} & \textbf{Purpose} \\
\midrule
Public Subnet 1 & 10.0.1.0/24 & us-east-1a & EMR master, Bastion host \\
Public Subnet 2 & 10.0.2.0/24 & us-east-1b & High availability \\
\bottomrule
\end{tabular}
\end{table}

\textbf{Justification:} Two public subnets in different AZs provide fault tolerance. EMR requires multiple AZs for high availability, and having subnets in different AZs is an EMR best practice.

\subsection{Internet Gateway \& Route Table}

\textbf{Internet Gateway:} \texttt{25tfgf-igw}
\begin{itemize}[nosep]
    \item Attached to the VPC to enable internet access for resources in public subnets
\end{itemize}

\textbf{Route Table:} \texttt{25tfgf-public-rt}
\begin{table}[H]
\centering
\begin{tabular}{@{}lll@{}}
\toprule
\textbf{Destination} & \textbf{Target} & \textbf{Purpose} \\
\midrule
10.0.0.0/16 & Local & Internal VPC traffic \\
0.0.0.0/0 & igw-xxxxxx & All internet traffic via IGW \\
\bottomrule
\end{tabular}
\end{table}

\subsection{Security Group Rules}

\textbf{EC2 Security Group (}\texttt{25tfgf-ec2-sg}\textbf{):}
\begin{table}[H]
\centering
\begin{tabular}{@{}lll@{}}
\toprule
\textbf{Port} & \textbf{Source} & \textbf{Purpose} \\
\midrule
22 & 0.0.0.0/0 & SSH access for administration \\
11434 & 0.0.0.0/0 & Ollama API (LLM runner) \\
8080 & 0.0.0.0/0 & OpenWebUI web interface \\
\bottomrule
\end{tabular}
\end{table}

\textbf{Justification:} SSH on port 22 is needed for server administration. Port 11434 (Ollama) and 8080 (OpenWebUI) must be accessible for the chatbot to function.

\textbf{EMR Master Security Group (}\texttt{25tfgf-emr-master-sg}\textbf{):}
\begin{table}[H]
\centering
\begin{tabular}{@{}lll@{}}
\toprule
\textbf{Port} & \textbf{Source} & \textbf{Purpose} \\
\midrule
22 & 0.0.0.0/0 & SSH access \\
9443 & 0.0.0.0/0 & JupyterHub web interface \\
0--65535 & 10.0.0.0/16 & EMR internal communication \\
\bottomrule
\end{tabular}
\end{table}

\textbf{EMR Slave Security Group (}\texttt{25tfgf-emr-slave-sg}\textbf{):}
\begin{table}[H]
\centering
\begin{tabular}{@{}lll@{}}
\toprule
\textbf{Port} & \textbf{Source} & \textbf{Purpose} \\
\midrule
22 & 0.0.0.0/0 & SSH access \\
0--65535 & 10.0.0.0/16 & EMR internal communication \\
\bottomrule
\end{tabular}
\end{table}

\textbf{Justification:} EMR security groups allow all internal communication (10.0.0.0/16) so the master can communicate with slave nodes. Port 9443 for JupyterHub is needed to access the EMR notebook interface during development.

\subsection{Terraform Files}

\textit{[CONFIRM: All Terraform files committed to GitHub in /terraform/ directory]}

\begin{lstlisting}[style=bash]
# Terraform files
terraform/
|-- main.tf            # VPC, subnets, IGW, route table
|-- variables.tf       # net_id, key_name, region
|-- outputs.tf         # VPC ID, subnet IDs, security group IDs
|-- provider.tf        # AWS provider configuration
|-- security-groups.tf # EC2 and EMR security groups
|-- emr.tf             # EMR IAM roles and instance profile
|-- ec2.tf             # EC2 instance configuration
\end{lstlisting}

\textbf{Justification for using Terraform over AWS Console:} Terraform was chosen because it provides Infrastructure-as-Code (IaC) capabilities, making the entire infrastructure reproducible with a single \texttt{terraform apply} command. This eliminates manual configuration errors, ensures consistency across deployments, and allows version-controlled infrastructure changes. Additionally, Terraform's state management simplifies tracking resource dependencies and teardown (\texttt{terraform destroy}), which is critical for a course project where cost control is essential.

\section{Model \& Dataset Selection (3 marks)}

\subsection{Model Selection: Llama-3.2-1B-Instruct (via Unsloth)}

\begin{table}[H]
\centering
\begin{tabular}{@{}ll@{}}
\toprule
\textbf{Attribute} & \textbf{Value} \\
\midrule
Model Name & unsloth/Llama-3.2-1B-Instruct \\
Base Model & meta-llama/Llama-3.2-1B-Instruct \\
Parameters & 1 billion \\
Source & \url{https://huggingface.co/unsloth/Llama-3.2-1B-Instruct} \\
License & Llama 3.2 License (requires acceptance on HuggingFace) \\
Quantization & 4-bit NF4 via QLoRA for fine-tuning \\
Format & GGUF (Q4\_K\_M) for Ollama deployment \\
\bottomrule
\end{tabular}
\end{table}

\textbf{Justification for Model Selection:}

\begin{enumerate}
    \item \textbf{Parameter Count:} At 1B parameters, the model is well under the 10B requirement, making it extremely suitable for local fine-tuning on an RTX 5000 Ada (16GB VRAM) with headroom to spare.
    \item \textbf{Hardware Requirements:} The 4-bit quantized version requires approximately 4--5GB VRAM with QLoRA, enabling fine-tuning comfortably on the RTX 5000 Ada workstation.
    \item \textbf{Domain Fit for Medical Q\&A:} Llama-3.2-1B-Instruct is a strong general-knowledge model with excellent instruction-following capabilities, making it well-suited for question-answering tasks. Its safety tuning helps it avoid generating harmful medical advice, while its broad pre-training corpus includes biomedical literature that provides a foundation for medical domain adaptation. The chat-template format natively supports multi-turn conversational interactions, which aligns with the symptom-to-advice workflow of a medical assistant.
    \item \textbf{Unsloth Optimization:} Using \texttt{unsloth/Llama-3.2-1B-Instruct} provides 2x faster training and 70\% less VRAM usage compared to the standard HuggingFace implementation.
    \item \textbf{Instruction Tuned:} The \texttt{-Instruct} variant is already instruction-tuned, providing a better foundation for medical domain adaptation through QLoRA fine-tuning.
    \item \textbf{Alternative Considered:} \texttt{unsloth/gemma-2-2b-it} was considered, but the Llama-3.2-1B-Instruct model offers similar quality with half the parameters, leading to faster training and lower inference latency on the CPU-only EC2 instance.
\end{enumerate}

\subsection{Dataset Selection: AI Medical Dataset}

\begin{table}[H]
\centering
\begin{tabular}{@{}ll@{}}
\toprule
\textbf{Attribute} & \textbf{Value} \\
\midrule
Dataset Name & ruslanmv/ai-medical-dataset \\
Source & \url{https://huggingface.co/datasets/ruslanmv/ai-medical-dataset} \\
License & CC-BY 4.0 \\
Total Samples & 21,210,000 (21.2M) \\
Original Split & Train only (provider-derived, no official split) \\
Columns & \texttt{question} (medical question), \texttt{context} (clinical context) \\
\bottomrule
\end{tabular}
\end{table}

\textbf{Data Sources:}
\begin{table}[H]
\centering
\begin{tabular}{@{}lr@{}}
\toprule
\textbf{Source} & \textbf{Word Count} \\
\midrule
ClinicalTrials & 127.4M words \\
EMEA & 12M words \\
PubMed & 968.4M words \\
\bottomrule
\end{tabular}
\end{table}

\subsection{Train/Validation/Test Split Strategy}

After Spark preprocessing on EMR, the dataset is split:

\begin{table}[H]
\centering
\begin{tabular}{@{}llll@{}}
\toprule
\textbf{Split} & \textbf{Ratio} & \textbf{Sample Size} & [TO\_UPDATE: Actual counts] \\
\midrule
Train & 80\% & 50,000 & [CONFIRM] \\
Validation & 10\% & 2,500 & [CONFIRM] \\
Test & 10\% & --- & [CONFIRM] \\
\bottomrule
\end{tabular}
\end{table}

\textbf{Actual Filtered Dataset Size:} The raw dataset contains 21.2M records. After applying the Spark preprocessing filters (context length: 200--4096 chars, question length: $\geq$20 chars), approximately [TO\_UPDATE: X] records remained. From this filtered set, 50,000 training samples and 2,500 evaluation samples were randomly selected for fine-tuning. The exact post-filter count can be verified from the Spark EDA output.

\textbf{Data Leakage Prevention:} The split uses a random 80/10/10 distribution via Spark's \texttt{randomSplit()}. Since this is a self-supervised fine-tuning task (not a held-out challenge), no external test set is required. The random split ensures no data from the test portion is used during training or hyperparameter tuning.

\subsection{Sample Data (Verbatim)}

\textbf{Example from dataset:}
\begin{lstlisting}
Question: What is the mechanism of action of ibuprofen?
Context: Ibuprofen is a nonsteroidal anti-inflammatory drug (NSAID) that works by
inhibiting the cyclooxygenase (COX) enzymes, which are responsible for producing
prostaglandins. By blocking COX-1 and COX-2, ibuprofen reduces inflammation, pain,
and fever by decreasing the production of prostaglandins that sensitize pain
receptors and promote inflammation.
\end{lstlisting}

\textbf{After preprocessing (chat template applied in fine-tuning script):}

The fine-tuning script (\texttt{colab/fine\_tune.py}) reformats each example into a multi-turn chat conversation using the model's native chat template (\texttt{tokenizer.apply\_chat\_template}).

The script performs two key text-processing steps before templating:
\begin{enumerate}
    \item \textbf{Relevant-sentence extraction} (\texttt{extract\_relevant\_sentence}): Scans the clinical context for the sentence(s) most semantically related to the question (using keyword overlap and medical-keyword boosting), then returns the best-matching sentence plus the next two sentences, capped at $\sim$600 characters.
    \item \textbf{Second-person rewriting} (\texttt{rewrite\_to\_second\_person}): Converts third-person clinical prose (e.g. ``The patient has been experiencing...'') into direct second-person advice (e.g. ``You have been experiencing...''). A 30-pattern regex pipeline handles verb-agreement fixes and removes academic filler phrases (``We show that...'', ``In conclusion,...'', etc.).
\end{enumerate}

A system prompt enforces the response style:
\begin{lstlisting}
You are a helpful medical assistant. The user describes their symptoms or asks a medical question.
Respond directly to THEM using 'you' and 'your'. Be concise (1-3 sentences).
Do NOT describe patients in the third person. Do NOT use clinical note style.
Do NOT use bullet points or lists. Write like you're talking to the person asking.
\end{lstlisting}

\textbf{Resulting chat-template example:}
\begin{lstlisting}
<|start_header_id|>system<|end_header_id|>
You are a helpful medical assistant...<|eot_id|>
<|start_header_id|>user<|end_header_id|>
What is the mechanism of action of ibuprofen?<|eot_id|>
<|start_header_id|>assistant<|end_header_id|>
You can take ibuprofen to reduce inflammation and pain because it blocks COX enzymes...
\end{lstlisting}

\subsection{Summary Statistics}

\textit{[TO\_UPDATE: Insert EDA figures and statistics from your Spark EDA run]}

\begin{table}[H]
\centering
\begin{tabular}{@{}ll@{}}
\toprule
\textbf{Statistic} & \textbf{Value} \\
\midrule
Mean token length (train) & [CONFIRM] \\
Median token length (train) & [CONFIRM] \\
95th percentile token length & [CONFIRM] \\
Min context length & 200 chars (filter threshold) \\
Max context length & 4096 chars (filter threshold) \\
\bottomrule
\end{tabular}
\end{table}

\textbf{Note on Label/Class Balance:} This is a text-generation (not classification) dataset; there are no discrete labels or classes. Consequently, traditional class-balance analysis is not applicable. Instead, the EDA focuses on length distributions (token count, context length, question length) and split proportions, which are the relevant quality metrics for generative fine-tuning.

\textit{[INSERT: Token Length Distribution Figure with caption]}\\
\textit{Figure 1: Token length distribution of the processed training dataset. The distribution is right-skewed with most samples between 200--1500 tokens.}

\textit{[INSERT: Context Length Distribution Figure with caption]}\\
\textit{Figure 2: Context length distribution showing the filtering preserved samples within the 200--4096 character range.}

\textit{[INSERT: Question Length Distribution Figure with caption]}\\
\textit{Figure 3: Question length distribution after the minimum 20-character filter was applied.}

\section{Data Preprocessing with Apache Spark on EMR (5 marks)}

\subsection{EMR Cluster Configuration}

\textit{[INSERT: EMR console screenshot showing cluster configuration]}

\begin{table}[H]
\centering
\begin{tabular}{@{}ll@{}}
\toprule
\textbf{Component} & \textbf{Value} \\
\midrule
Cluster Name & 25tfgf-emr-cluster \\
Release Label & emr-7.2.0 \\
Region & us-east-1 \\
Master Node & m5.xlarge (4 vCPU, 16GB RAM) \\
Core Nodes & 2 $\times$ m5.xlarge (4 vCPU, 16GB RAM each) \\
Total Nodes & 3 \\
Applications & Spark 3.5.0, JupyterHub \\
Bootstrap & Custom script (s3://25tfgf-ai-medical/bootstrap\_emr.sh) \\
\bottomrule
\end{tabular}
\end{table}

\textbf{Instance Type Justification:} m5.xlarge provides a good balance of cost and performance for Spark processing. The 16GB RAM allows efficient processing of the 100k sample dataset. Using m5 (not m4) provides better price-performance due to newer hardware.

\subsection{EMR Bootstrap Script}

The EMR cluster uses a custom bootstrap script (\texttt{spark/bootstrap\_emr.sh}) uploaded to S3 at \texttt{s3://25tfgf-ai-medical/bootstrap\_emr.sh}. This script runs on every node during cluster initialization and performs the following:

\begin{enumerate}
    \item \textbf{System updates:} Runs \texttt{yum update} to ensure the latest Amazon Linux packages.
    \item \textbf{Python dependencies:} Installs required Python packages from \texttt{spark/requirements.txt} (e.g., \texttt{pyarrow}, \texttt{pandas}, \texttt{boto3}) into the EMR Python environment so the PySpark pipeline can execute successfully.
    \item \textbf{Directory setup:} Creates necessary local directories for Spark temp files and logs.
\end{enumerate}

\textit{[INSERT: Full bootstrap\_emr.sh contents or a screenshot of the S3 object showing the script]}

\subsection{PySpark Pipeline (\texttt{spark/preprocess.py})}

\textit{[CONFIRM: Full script committed to GitHub at /spark/preprocess.py]}

The preprocessing pipeline performs the following steps:

\begin{enumerate}
    \item \textbf{Load:} Read raw Parquet files from S3 (\texttt{s3://25tfgf-ai-medical/raw/})
    \item \textbf{Filter:} Apply quality filters:
    \begin{itemize}[nosep]
        \item Context length: 200--4096 characters
        \item Question length: $\geq$20 characters
    \end{itemize}
    \item \textbf{Parse:} Extract and clean \texttt{question\_parsed} and \texttt{context\_parsed} columns
    \item \textbf{Split:} Random 80/10/10 train/val/test split using Spark's \texttt{randomSplit()}
    \item \textbf{Write:} Save as partitioned Parquet to S3 (\texttt{s3://25tfgf-ai-medical/processed/})
\end{enumerate}

\textbf{Note:} The fine-tuning script (\texttt{colab/fine\_tune.py}) applies the chat template and second-person rewriting at training time (see Section 3), so the Spark pipeline outputs raw parsed question/context pairs rather than a fixed prompt string.

\subsection{Preprocessed Output Files}

\textit{[INSERT: S3 console screenshot showing processed parquet files in s3://25tfgf-ai-medical/processed/]}

Output structure:
\begin{lstlisting}
processed/
|-- _SUCCESS
|-- part-00000-*.parquet  (train partition)
|-- part-00001-*.parquet  (train partition)
|-- ...
`-- (validation and test partitions)
\end{lstlisting}

\subsection{EMR Cluster Teardown (REQUIRED)}

\textit{[INSERT: EMR console screenshot showing cluster in TERMINATED state]}

\textbf{CRITICAL:} The EMR cluster was terminated immediately after preprocessing to avoid ongoing charges. The cluster ran for approximately [TO\_UPDATE: duration] and cost approximately [TO\_UPDATE: cost].

\section{Model Fine-Tuning (6 marks)}

\subsection{Fine-Tuning Configuration}

\begin{table}[H]
\centering
\begin{tabular}{@{}ll@{}}
\toprule
\textbf{Component} & \textbf{Value} \\
\midrule
Library & Unsloth (local workstation + \texttt{unsloth/Llama-3.2-1B-Instruct} base model) \\
Technique & QLoRA (4-bit Normal Float 4 quantization + Low-Rank Adaptation) \\
Hardware & Local workstation (RTX 5000 Ada, 16 GB VRAM) \\
Sample Size & 50,000 train / 2,500 eval (from processed dataset) \\
Code & \texttt{/colab/fine\_tune.py} \\
\bottomrule
\end{tabular}
\end{table}

\subsection{Unsloth Benefits}

\begin{itemize}[nosep]
    \item \textbf{2--5x faster training} compared to standard HuggingFace implementation
    \item \textbf{70\% less VRAM usage} (4--5GB vs $\sim$16GB for full fine-tuning)
    \item \textbf{Pre-installed dependencies:} transformers, peft, trl, bitsandbytes
    \item \textbf{Native Llama 3.2 support} with optimized attention mechanisms
\end{itemize}

\subsection{Training Details}

The fine-tuning was performed on a local workstation (RTX 5000 Ada, 16 GB VRAM) using the \texttt{unsloth/Llama-3.2-1B-Instruct} model (Unsloth's optimized Llama 3.2 1B instruction-tuned variant) and the processed medical dataset from S3. The training script (\texttt{/colab/fine\_tune.py}) downloads the processed Parquet files from the \texttt{25tfgf-ai-medical} S3 bucket, applies the custom medical chat template with second-person rewriting (see Section 3), and runs QLoRA fine-tuning with the hyperparameters below.

\textbf{Why QLoRA over Full Fine-Tuning:} Full fine-tuning would require updating all 1B parameters simultaneously, demanding approximately 16GB+ VRAM (model weights + gradients + optimizer states) --- exceeding the RTX 5000 Ada's capacity. QLoRA addresses this by:
\begin{enumerate}
    \item \textbf{4-bit Quantization:} Base model weights are quantized to 4-bit NF4, reducing memory footprint by $\sim$4$\times$.
    \item \textbf{Low-Rank Adapters:} Only small rank-16 adapter matrices are trained ($\sim$0.1\% of total parameters), drastically reducing gradient and optimizer state memory.
    \item \textbf{Double Quantization:} Additional quantization of adapter gradients further minimizes VRAM usage.
\end{enumerate}

This combination allows fine-tuning on a 16GB consumer GPU while achieving comparable downstream performance to full fine-tuning. The resulting adapters are also lightweight ($\sim$10--20MB), making them easy to version-control, share, and deploy.

\textbf{Checkpoint Resumption \& Idempotency:} The script implements two robust mechanisms to prevent accidental re-training and support resumption:
\begin{enumerate}
    \item \textbf{Final-model guard:} Before training starts, \texttt{final\_model\_exists(FINAL\_DIR)} checks for \texttt{adapter\_config.json} and \texttt{adapter\_model.safetensors} (or \texttt{.bin}). If the final LoRA adapter is already present, \texttt{SKIP\_TRAIN} is set to \texttt{True} and training is bypassed entirely.
    \item \textbf{Checkpoint resumption:} If training is interrupted, \texttt{get\_last\_checkpoint(CHECKPOINT\_DIR)} scans the checkpoint directory for the highest-numbered \texttt{checkpoint-*} folder. The \texttt{trainer.train(resume\_from\_checkpoint=...)} call automatically resumes from that checkpoint, preserving optimizer state and step count. This is essential for long training runs on limited GPU time.
\end{enumerate}

\subsection{Hyperparameter Table}

\begin{longtable}{@{}p{4.2cm}p{2.8cm}p{6cm}@{}}
\toprule
\textbf{Parameter} & \textbf{Value} & \textbf{Rationale} \\
\midrule
\endhead
Learning Rate & 2e-5 & Conservative rate for stable medical-domain adaptation \\
Batch Size (per device) & 4 & Fits comfortably in RTX 5000 Ada VRAM with 4-bit weights \\
Gradient Accumulation & 4 & Effective batch size of 16 \\
Epochs & 1 & Time constraint \\
LoRA Rank (r) & 16 & Sufficient rank for a 1B-parameter model \\
LoRA Alpha & 16 & 1$\times$ rank scaling factor \\
LoRA Dropout & 0 & No dropout (standard for LoRA fine-tuning) \\
LoRA Target Modules & q\_proj, k\_proj, v\_proj, o\_proj, gate\_proj, up\_proj, down\_proj & All linear projection layers adapted \\
Max Sequence Length & 1024 (training) / 512 (model load) & Truncates longer sequences to fit GPU memory \\
Quantization & 4-bit NF4 & Optimal for QLoRA \\
Optimizer & adamw\_8bit & Memory-efficient AdamW from Unsloth \\
Warmup Ratio & 0.03 & $\sim$3\% of total steps for stable early training \\
LR Scheduler & cosine & Cosine decay for stable convergence \\
Weight Decay & 0.01 & Light regularization \\
Max Grad Norm & 1.0 & Gradient clipping for stability \\
Logging Steps & 50 & Frequent loss logging \\
Eval Steps & 500 & Periodic validation evaluation \\
Save Steps & 1000 & Checkpointing with \texttt{save\_total\_limit=2} \\
\bottomrule
\end{longtable}

\subsection{Training Code}

\textit{[CONFIRM: Runnable notebook committed to GitHub at /colab/fine\_tune.ipynb and training script at /colab/fine\_tune.py]}

The training script (\texttt{/colab/fine\_tune.py}) downloads processed data from S3 and fine-tunes using Unsloth:

\begin{lstlisting}
from unsloth import FastLanguageModel
import torch

# Load Llama-3.2-1B-Instruct via Unsloth
model, tokenizer = FastLanguageModel.from_pretrained(
    "unsloth/Llama-3.2-1B-Instruct",
    max_seq_length     = 512,
    dtype              = None,
    load_in_4bit       = True,
    device_map         = "cuda",
)

# Configure LoRA adapters
model = FastLanguageModel.get_peft_model(
    model,
    r                = 16,
    target_modules   = ["q_proj", "k_proj", "v_proj", "o_proj",
                        "gate_proj", "up_proj", "down_proj"],
    lora_alpha       = 16,
    lora_dropout     = 0,
    bias             = "none",
    use_gradient_checkpointing = "unsloth",
    random_state     = 3407,
)

# Load dataset from S3 (50k train / 2.5k eval)
# ... (see /colab/fine_tune.py for full implementation)

# Train with SFTTrainer
trainer = SFTTrainer(
    model=model,
    tokenizer=tokenizer,
    train_dataset=train_data,
    eval_dataset=eval_data,
    dataset_text_field="text",
    max_seq_length=1024,
    dataset_num_proc=2,
    packing=True,
    args=TrainingArguments(
        per_device_train_batch_size=4,
        gradient_accumulation_steps=4,
        num_train_epochs=1,
        max_steps=-1,
        learning_rate=2e-5,
        warmup_ratio=0.03,
        lr_scheduler_type="cosine",
        weight_decay=0.01,
        max_grad_norm=1.0,
        fp16=not is_bfloat16_supported(),
        bf16=is_bfloat16_supported(),
        logging_steps=50,
        eval_strategy="steps",
        eval_steps=500,
        save_strategy="steps",
        save_steps=1000,
        save_total_limit=2,
        optim="adamw_8bit",
        seed=3407,
        output_dir=str(CHECKPOINT_DIR),
        report_to="none",
    ),
)
trainer.train()

# Save LoRA adapter + tokenizer
trainer.save_model("./cloud_project/final_lora")
tokenizer.save_pretrained("./cloud_project/final_lora")

# Export GGUF for Ollama
model.save_pretrained_gguf(
    "./cloud_project/medical_assistant_gguf",
    tokenizer,
    quantization_method="q4_k_m"
)
\end{lstlisting}

\subsection{Base vs Fine-Tuned Model Comparison}

\textbf{Model:} \texttt{unsloth/Llama-3.2-1B-Instruct} fine-tuned as \texttt{medical-assistant}

\textbf{Example 1:}

\textit{Prompt:} ``I am a 45-year-old male. I have been experiencing severe chest pain radiating to my left arm, accompanied by shortness of breath and sweating for the past 2 hours. What could be causing these symptoms and how urgent is this?''

\textit{Base Model Response (unsloth/Llama-3.2-1B-Instruct without fine-tuning):}\\
\hspace*{1em}\textit{[TO\_UPDATE: Insert base model response]}

\textit{Fine-Tuned Model Response (medical-assistant):}\\
\hspace*{1em}\textit{[TO\_UPDATE: Insert fine-tuned model response]}

\vspace{0.5em}
\textbf{Example 2:}

\textit{Prompt:} ``My doctor prescribed me Metformin 500mg twice daily for Type 2 Diabetes. What are the common and serious side effects I should watch out for?''

\textit{Base Model Response:}\\
\hspace*{1em}\textit{[TO\_UPDATE: Insert base model response]}

\textit{Fine-Tuned Model Response:}\\
\hspace*{1em}\textit{[TO\_UPDATE: Insert fine-tuned model response]}

\vspace{0.5em}
\textbf{Example 3:}

\textit{Prompt:} ``I am a 32-year-old non-smoker with no prior lung conditions. I have had a dry, persistent cough for 3 weeks with mild fatigue but no fever. What are the possible causes and when should I seek medical attention?''

\textit{Base Model Response:}\\
\hspace*{1em}\textit{[TO\_UPDATE: Insert base model response]}

\textit{Fine-Tuned Model Response:}\\
\hspace*{1em}\textit{[TO\_UPDATE: Insert fine-tuned model response]}

\subsection{Qualitative Analysis}

\textit{[TO\_UPDATE: Provide analysis comparing responses, highlighting improvements in medical accuracy, domain terminology, and response structure]}

\subsection{Optional: Training Loss Curve}

\textit{[INSERT: Training loss curve figure if available from notebook output]}

\subsection{Model Upload to S3 (Bridge to Deployment)}

After fine-tuning and GGUF export are complete on the local workstation, the model artifacts must be transferred to S3 so the EC2 instance can download them for serving.

\textbf{Artifacts produced:}
\begin{table}[H]
\centering
\begin{tabular}{@{}lll@{}}
\toprule
\textbf{Artifact} & \textbf{Local Path} & \textbf{S3 Destination} \\
\midrule
LoRA adapter & \texttt{./cloud\_project/final\_lora/} & \texttt{s3://25tfgf-ai-medical/models/final\_lora/} \\
GGUF model & \texttt{./cloud\_project/medical\_assistant\_gguf/} & \texttt{s3://25tfgf-ai-medical/models/medical\_assistant\_gguf/} \\
Training logs & \texttt{./cloud\_project/training\_log.csv} & \texttt{s3://25tfgf-ai-medical/models/training\_log.csv} \\
Loss curve & \texttt{./cloud\_project/training\_loss\_curve.png} & \texttt{s3://25tfgf-ai-medical/models/training\_loss\_curve.png} \\
\bottomrule
\end{tabular}
\end{table}

\textbf{Upload commands (run from local workstation):}
\begin{lstlisting}[style=bash]
aws s3 sync ./cloud_project/final_lora/ s3://25tfgf-ai-medical/models/final_lora/
aws s3 sync ./cloud_project/medical_assistant_gguf/ s3://25tfgf-ai-medical/models/medical_assistant_gguf/
aws s3 cp ./cloud_project/training_log.csv s3://25tfgf-ai-medical/models/
aws s3 cp ./cloud_project/training_loss_curve.png s3://25tfgf-ai-medical/models/
\end{lstlisting}

\textit{[INSERT: S3 console screenshot showing the uploaded model folders in s3://25tfgf-ai-medical/models/]}

\section{Model Deployment on EC2 (3 marks)}

\subsection{EC2 Instance Configuration}

\begin{table}[H]
\centering
\begin{tabular}{@{}ll@{}}
\toprule
\textbf{Component} & \textbf{Value} \\
\midrule
Instance Type & m5.xlarge \\
Region & us-east-1 \\
AMI & Ubuntu 22.04 LTS (ami-xxxxxxxx) \\
vCPUs & 4 \\
RAM & 16 GB \\
Storage & [TO\_UPDATE: EBS size] GB \\
\bottomrule
\end{tabular}
\end{table}

\textbf{Justification:} m5.xlarge provides sufficient CPU resources for serving the fine-tuned model via Ollama and hosting the OpenWebUI interface. The model is already 4-bit quantized and runs efficiently on CPU. This instance type offers a lower cost alternative to GPU instances for inference workloads.

\subsection{Ollama Installation Commands}

\textit{[CONFIRM: Commands copy-pasted verbatim in README.md]}

\begin{lstlisting}[style=bash]
# SSH to EC2
ssh -i 25tfgf-key.pem ubuntu@<EC2_PUBLIC_IP>

# Update system
sudo apt-get update && sudo apt-get upgrade -y

# Install Ollama
curl -fsSL https://ollama.com/install.sh | sh

# Create systemd service for auto-start
sudo nano /etc/systemd/system/ollama.service
# [Insert ollama.service content -- see below]

# Enable and start
sudo systemctl daemon-reload
sudo systemctl enable ollama
sudo systemctl start ollama

# Verify
sudo systemctl status ollama --no-pager
\end{lstlisting}

\subsection{Ollama systemd Service (\texttt{ollama.service})}

The following systemd service file ensures Ollama starts automatically on boot and restarts on failure:

\begin{lstlisting}[style=ini]
[Unit]
Description=Ollama Service
After=network-online.target

[Service]
ExecStart=/usr/local/bin/ollama serve
User=ubuntu
Group=ubuntu
Restart=always
RestartSec=3
Environment="PATH=/usr/local/sbin:/usr/local/bin:/usr/sbin:/usr/bin:/sbin:/bin"

[Install]
WantedBy=default.target
\end{lstlisting}

\textbf{Justification:} \texttt{Restart=always} ensures the Ollama server recovers from crashes or unexpected terminations. \texttt{After=network-online.target} guarantees network availability before starting, preventing race conditions during EC2 boot.

\subsection{Model Loading and Serving}

\begin{lstlisting}[style=bash]
# Download GGUF from S3 (after uploading from local workstation)
aws s3 cp s3://25tfgf-ai-medical/models/medical_assistant_gguf /home/ubuntu/model/ --recursive

# Create Ollama model from GGUF
ollama create medical-assistant -f /home/ubuntu/model/Modelfile

# Verify model is loaded
ollama list
\end{lstlisting}

\subsection{Ollama Modelfile}

The \texttt{Modelfile} tells Ollama how to load the GGUF weights and configures generation parameters and the system prompt:

\begin{lstlisting}[style=docker]
FROM /home/ubuntu/model/unsloth.Q4_K_M.gguf

PARAMETER temperature 0.4
PARAMETER top_p 0.85
PARAMETER repeat_penalty 1.15
PARAMETER num_ctx 2048

SYSTEM """You are a helpful medical assistant. The user describes their symptoms or asks a medical question. Respond directly to THEM using 'you' and 'your'. Be concise (1-3 sentences). Do NOT describe patients in the third person. Do NOT use clinical note style. Do NOT use bullet points or lists. Write like you're talking to the person asking."""
\end{lstlisting}

\textbf{Key parameters:}
\begin{itemize}[nosep]
    \item \texttt{temperature 0.4}: Low temperature for deterministic, factual medical responses.
    \item \texttt{top\_p 0.85}: Nucleus sampling to maintain coherence while allowing slight variation.
    \item \texttt{repeat\_penalty 1.15}: Discourages repetitive phrasing in longer responses.
    \item \texttt{num\_ctx 2048}: Sufficient context window for multi-turn medical conversations.
\end{itemize}

\subsection{Ollama Terminal Screenshot}

\textit{[INSERT: Terminal screenshot showing Ollama running with medical-assistant visible in model list]}

\subsection{API Test with curl}

\textit{[INSERT: Screenshot of curl command and response]}

\begin{lstlisting}[style=bash]
curl http://localhost:11434/api/generate -d '{
  "model": "medical-assistant",
  "prompt": "I have a severe headache, sensitivity to light, and nausea. What could this be?",
  "stream": false
}'
\end{lstlisting}

\section{Web Interface (2 marks)}

\subsection{OpenWebUI Configuration}

OpenWebUI is configured to start automatically via systemd service and connects to the local Ollama instance at \url{http://localhost:11434}.

\subsection{OpenWebUI systemd Service (\texttt{openwebui.service})}

The following systemd service file ensures OpenWebUI starts automatically on boot and restarts on failure:

\begin{lstlisting}[style=ini]
[Unit]
Description=OpenWebUI Service
After=network-online.target ollama.service
Requires=ollama.service

[Service]
Type=simple
ExecStart=/usr/bin/docker run --rm --name openwebui -p 8080:8080 \
  -e OLLAMA_API_BASE_URL=http://host.docker.internal:11434/api \
  -v open-webui:/app/backend/data \
  ghcr.io/open-webui/open-webui:main
ExecStop=/usr/bin/docker stop openwebui
ExecStopPost=/usr/bin/docker rm openwebui
User=ubuntu
Group=ubuntu
Restart=always
RestartSec=5

[Install]
WantedBy=default.target
\end{lstlisting}

\textbf{Justification:} \texttt{After=ollama.service} and \texttt{Requires=ollama.service} establish a dependency chain so OpenWebUI only starts after Ollama is ready. \texttt{Restart=always} with \texttt{RestartSec=5} ensures recovery from Docker container crashes. Using Docker simplifies deployment and dependency management.

\subsection{Web Interface Screenshot}

\textit{[INSERT: Browser screenshot of OpenWebUI running at http://<EC2\_PUBLIC\_IP>:8080 with model name medical-assistant visible in model selector]}

\subsection{Sample Conversation Screenshot}

\textit{[INSERT: Screenshot showing a complete sample conversation through the OpenWebUI interface]}

\section{GitHub Repository Requirements}

\textit{[CONFIRM: All requirements met and links verified]}

\begin{table}[H]
\centering
\begin{tabular}{@{}lll@{}}
\toprule
\textbf{Requirement} & \textbf{Status} & \textbf{Location} \\
\midrule
PySpark script & [CONFIRM] & /spark/preprocess.py \\
Fine-tuning notebook & [CONFIRM] & /colab/fine\_tune.ipynb \\
Fine-tuning script & [CONFIRM] & /colab/fine\_tune.py \\
Terraform files & [CONFIRM] & /terraform/*.tf \\
README.md with replication steps & [CONFIRM] & /README.md \\
Prerequisites documented & [CONFIRM] & /README.md \\
Cost summary table & [CONFIRM] & /README.md \\
\bottomrule
\end{tabular}
\end{table}

\section{Cost Summary}

\begin{table}[H]
\centering
\begin{tabular}{@{}lllr@{}}
\toprule
\textbf{Service} & \textbf{Configuration} & \textbf{Duration} & \textbf{Actual Cost} \\
\midrule
S3 & $\sim$5GB storage & Monthly & \$[TO\_UPDATE] \\
EMR & m5.xlarge (1+2 nodes) & $\sim$[TO\_UPDATE] minutes & \$[TO\_UPDATE] \\
EC2 & m5.xlarge & [TO\_UPDATE] hours & \$[TO\_UPDATE] \\
\midrule
\textbf{Total} & & & \textbf{\$[TO\_UPDATE]} \\
\bottomrule
\end{tabular}
\end{table}

\section{Mark Summary}

\begin{table}[H]
\centering
\begin{tabular}{@{}llcr@{}}
\toprule
\textbf{Section} & \textbf{Deliverable} & \textbf{Marks} & \textbf{Status} \\
\midrule
1 & System Architecture Diagram + Paragraph & 2/2 & [TO\_UPDATE] \\
2 & VPC \& Networking (Terraform) & 4/4 & [TO\_UPDATE] \\
3 & Model \& Dataset Selection & 3/3 & [TO\_UPDATE] \\
4 & EMR + Spark Preprocessing & 5/5 & [TO\_UPDATE] \\
5 & Model Fine-Tuning & 6/6 & [TO\_UPDATE] \\
6 & EC2 Deployment & 3/3 & [TO\_UPDATE] \\
7 & Web Interface & 2/2 & [TO\_UPDATE] \\
\midrule
\textbf{Total} & & \textbf{25/25} & \\
\bottomrule
\end{tabular}
\end{table}

\vfill
\begin{center}
\textit{Report prepared for CISC 886 -- Cloud Computing, Queen's University}
\end{center}

\end{document}
